(1) 在 $z$ 自旋基中 $|+x\rangle=(|+z\rangle+|-z\rangle)/\sqrt2$、$|-x\rangle=(|+z\rangle-|-z\rangle)/\sqrt2$，所以初态 $|+z\rangle$ 对两个 $x$ 结果的投影模平方均为 $1/2$。理想投影测量后，选择 $+x$ 或 $-x$ 的子样本分别处于相应本征态。
(2) 给定第一测量结果 $+x$，测后态为 $|+x\rangle$。再次测量 $\sigma_z$ 时，$P(+z|+x)=|\langle+z|+x\rangle|^2=1/2$；同样 $P(-z|+x)=1/2$，因为 $|+x\rangle$ 没有确定的 $z$ 自旋。
(3) 若不读取 $x$ 测量结果，非选择性投影过程产生 $\rho=\frac12|+x\rangle\langle+x|+\frac12|-x\rangle\langle-x|=I/2$。再测 $z$ 得 $P(+z)=\operatorname{tr}(\rho|+z\rangle\langle+z|)=1/2$，与按全概率公式 $\sum_{\pm x}P(+z|\pm x)P(\pm x)$ 一致。
(4) 若直接对未受中间干扰的初态 $|+z\rangle$ 再次测量 $\sigma_z$，必得 $+1$，概率为一。中间的 $x$ 投影把原有 $z$ 相干性抹去，产生混合态，因此最终概率降至一半；这不是仪器误差，而是非对易可观测量的测量反作用。
